% results_rq2

\section{RQ2: Accuracy Versus Interface}
\label{sec:rq2}

A second common framing in the review-mining and mobile-HCI literature treats
accuracy and interface polish as a trade-off users weigh against one another.
RQ2 asks whether devotional-app users punish inaccuracy more harshly than a
poor interface. As a class, they do not. That comparison, however, is the
wrong one to draw: the class-level average hides a severity gradient inside
accuracy itself, and the rating signal developers typically optimise against
conceals the worst of it.

\subsection{As Classes, Accuracy and Interface Cost the Same}

Using M1, a review-level mixed-effects model of star rating against complaint
aspect (\(n = 324{,}687\); all coefficients discussed in this section survive
Benjamini--Hochberg correction at \(q < 0.05\)), the accuracy family ---
prayer-time accuracy, qibla, madhab, and calculation method --- averages
\(-0.647\) stars, against \(-0.600\) for \texttt{ui\_design}. The gap is under
a twentieth of a star. The intuition that users punish a wrong qibla or prayer
time more harshly than a cluttered screen is not supported at the level of the
complaint class, and we do not claim accuracy wins as a class.

\subsection{The Class Average Conceals a Severity Gradient}

Table~\ref{tab:rq2-coefficients} and Figure~\ref{fig:rq2-coefficients} give the
full spread underlying that average. Accuracy is not one thing: prayer-time
accuracy alone costs \(-1.112\) stars (95\% CI \([-1.142, -1.083]\)), 1.85
times the interface penalty, with a confidence interval disjoint from
\texttt{ui\_design}'s. The family mean is pulled down by \texttt{madhab}
(\(-0.324\)) and \texttt{qibla} (\(-0.550\)), which are both rarer and milder
complaints. Getting the prayer time wrong is among the most costly failures in
the entire corpus; getting the madhab option wrong is among the least ---
and \texttt{madhab} additionally reverses sign in an ordinal robustness check
(\(\beta = +0.216\), \(p = 0.649\)), so no directional claim rests on it
either way.

\begin{quote}
\itshape ``this app is not accurate, prayers time are 1 hour ahead, please fix
it''
\end{quote}
\hspace*{1em}--- Athan: Prayer Times \& Al Quran, 1 star

\begin{table}[t]
  \centering
  \small
  \caption{M1 coefficients on star rating, all complaint aspects surviving BH-FDR at \(q < 0.05\) (\(n = 324{,}687\)).}
  \label{tab:rq2-coefficients}
  \begin{tabular}{lrr}
    \toprule
    Complaint & \(\beta\) & 95\% CI \\
    \midrule
    Intrusive ads            & \(-1.439\) & \([-1.460, -1.418]\) \\
    Prayer-time accuracy     & \(-1.112\) & \([-1.142, -1.083]\) \\
    Stability / bugs         & \(-0.938\) & \([-0.965, -0.910]\) \\
    Reminders / adhan        & \(-0.845\) & \([-0.868, -0.823]\) \\
    UI design                & \(-0.600\) & \([-0.664, -0.537]\) \\
    Calculation method       & \(-0.600\) & \([-0.769, -0.430]\) \\
    Qibla                    & \(-0.550\) & \([-0.594, -0.507]\) \\
    Madhab                   & \(-0.324\) & \([-0.526, -0.122]\) \\
    \bottomrule
  \end{tabular}
\end{table}

\begin{figure}[t]
  \centering
  \includegraphics[width=\linewidth]{fig05_rq2_complaint_effects}
  \caption{M1 coefficients by complaint aspect, colour-coded by family. The spread within the accuracy family is visible directly in the error bars.}
  \label{fig:rq2-coefficients}
\end{figure}

\subsection{Qibla Failure Is Instability, Not Miscalculation}

The larger of the two qibla-related topics in our topic model (1,771 reviews)
is labelled around wrong direction, outweighing a second, smaller topic
(698) labelled around an inaccurate compass pointer. Reading the representative
reviews in the larger topic, the complaint is not that the app computes the
wrong bearing: users describe the compass returning a different bearing from
the same physical spot as the phone is rotated. We report this as a sensor and
calibration-handling problem rather than an error in the underlying qibla
computation, since the two failure modes call for different fixes.

\subsection{Star Ratings Hide Accuracy Defects}

The strongest result in RQ2 is methodological rather than substantive. Of
7,942 reviews carrying at least one negative accuracy-family complaint, 3,306
(41.6\%) awarded four or five stars, and 1,830 (23.0\%) awarded the maximum of
five; the full distribution runs 2,319 one-star, 914 two-star, 1,403
three-star, 1,476 four-star, and 1,830 five-star. Users report that the prayer
time or qibla direction is wrong and rate the application highly in the same
review.
Any prioritisation process that reads defect severity off the star
distribution --- which is to say, the signal most developers and app-store
ranking systems optimise against --- will under-weight accuracy defects by a
wide margin, because a substantial share of the reviews that name them never
produce a low rating.

\subsection{Advertising Costs More Than Any Accuracy Complaint}

One coefficient in Table~\ref{tab:rq2-coefficients} falls outside the
accuracy/interface contrast entirely and is worth flagging here: intrusive
advertising costs \(-1.439\) stars, above every accuracy-family term,
including prayer-time accuracy. We return to advertising as a recurring cost
across the bloat and gap analyses in \S\ref{sec:rq5} and \S\ref{sec:rq6}.

\subsection{Summary}

Accuracy and interface complaints are indistinguishable as classes, but that
comparison obscures more than it reveals. Prayer-time accuracy alone is a
severe, frequent failure whose cost is masked twice: once by averaging it with
milder accuracy complaints such as madhab and qibla, and again by arriving
inside four- and five-star reviews at a rate of 41.6\%. Star ratings are an
unsafe proxy for defect severity in this domain.
